\chapter{Additional info for high-profile individuals}
\label{ch:high-profile}

If you are a high-profile target, you'll evidently realize
that some of these initial suggestions fall short within
your threat model. This section provides additional
recommendations for profiles like yours.

Web3 professionals who have elevated privileges or
profiles -- for example, access to multisig treasury keys,
leadership roles, or possession of sensitive
organizational secrets are high-profile targets. These
users may be targeted more deliberately by criminals or
even nation-states. In addition to all the precautions
above, high-risk travelers should take further steps.

For high-profile or recognizable individuals, keeping a
low profile is essential. Beyond avoiding branded
merchandise, a simple yet effective tactic is wearing an
N95 mask or similar face covering. It's socially accepted,
draws no attention, and helps protect your identity --
making it harder for adversaries to target or track you
during public events.

\section{Use loaner or ``burner'' devices}

If feasible, travel with clean devices that don't contain
sensitive data. Leave your primary laptop/phone at a
protected location (assuming you also have security in
place as well), and bring a wiped, minimal laptop or a
cheap travel phone with only the basics. Log into what you
absolutely need (through secure channels) and nothing
more. Treat these devices as expendable -- assume they
will be compromised and plan to wipe them afterward. For
example, a senior developer might bring a laptop with no
source code on it and use a VPN or VDI (virtual desktop)
to access company systems when necessary, leaving no data
on the local disk. A hardware wallet keyholder might carry
a secondary hardware wallet with lower privileges (or a
single key to a multisig instead of full access).

Keep your primary keys secured in a location that is not
accompanying you. Remember, a nearly empty device can
raise suspicion at some borders, so don't make it obvious
-- load some innocuous data (music, generic files) so it
looks used, but nothing that would be harmful if
inspected.

\section{Plan for customs and border checks}

High-risk individuals may face increased scrutiny or
device searches at borders due to the sensitive nature of
their data or their roles (e.g., journalists). Before
crossing borders, purge beyond recovery or secure
sensitive information. Turn off devices before landing.
Some experts even recommend encrypting and then powering
down devices -- a powered-off device with strong
encryption is extremely hard to
access.\footnote{\url{https://www.reddit.com/r/technology/comments/1jl9dxg/how_to_lock_down_your_phone_if_youre_traveling_to}}

Disable cloud auto-sync of sensitive data; you don't want
customs inadvertently accessing company cloud drives if
they inspect your laptop. If asked to unlock devices,
having them powered off gives you an opportunity to state
that it's encrypted and requires a passphrase (which you
should have memorized, not written down).

It's wise for high-risk travelers to disable biometrics
entirely before travel -- use PIN/password only, as
mentioned -- so you cannot be compelled or tricked via
fingerprint/face.

Know your legal rights in the countries you transit; in
some places you can refuse to unlock (though it may mean
the device is held or you are denied entry), while in
others you might face penalties. This is a personal risk
decision -- but the best case is to carry nothing truly
incriminating or irreplaceable across a checkpoint.

For ultra-sensitive data, use secure communication
channels to retrieve it at destination rather than
carrying it. For instance, an executive could store
confidential files in a secure cloud drive they access
over VPN once abroad, rather than carrying them on a
laptop.

Border/immigration checks often require you to share your
travel plans, including hotel information and airline
details. Having a physical copy of them (printouts,
ingress-egress) means you do not need to access your
devices in front of agents or potential cameras while
being scrutinized. This will also prevent the need to let
them peek at them to provide information, thereby risking
the disclosure of sensitive personal information.

Since COVID, many nations have moved to digital
immigration paperwork. This often requires the use of an
email address. You should consider a dedicated email
address for interaction with government agencies, instead
of an email address that might have a public/open-source
affiliation with cryptocurrency (or other sensitive
topics).

\section{Enhanced device protections}

High-risk users should layer additional defenses. Lockdown
Mode (on iOS) or Android's equivalent secure modes should
be enabled if there's any chance of targeted spyware
(these modes disable exploit-prone services and
attachments). Use messaging apps with end-to-end
encryption and disappearing messages (NOT Telegram, Signal
is a good example) for any sensitive communications --
assume that standard SMS or emails could be monitored.

Consider using a Faraday bag for phones when not in use,
if you suspect active tracking or exploitation (this
prevents any signals in or out, though use sparingly as it
also blocks your calls).

If you leave a device in your hotel, you can put it in a
tamper-evident bag and seal it, or at least take measures
like noting the exact placement or taking a photo, so you
can detect if it was accessed.

Some high-risk individuals even weigh their devices before
and after travel to detect the addition of hardware
implants (a change in weight could indicate something like
a chip added) -- this is an extreme step, but shows the
level of caution possible. At the very least, physically
inspect your devices for new scratches, screws that look
tampered with, or unexpected behavior.

\section{Secure crypto keys with multi-party control}

If you have access to significant crypto funds (exchanges,
DAO treasury, etc.), implement policies that prevent a
single point of failure while you're traveling. For
example, if you're one of M-of-N multisig signers,
consider temporarily requiring an extra signer for
transactions while you're away (so if normally 2-of-3 can
move funds, bump it to 3-of-3 or add a 4th backup signer)
so that a compromised key or coerced action cannot alone
execute a transfer.

If you hold hardware keys, keep them geographically split
-- e.g.\ bring one key device with you, leave the backup
key in a safe place at home, and perhaps give a third key
to a colleague, so that even a forced disclosure cannot
result in an immediate loss of funds without
collaboration. Use duress codes if your hardware supports
it (some wallets allow a secondary PIN that opens a decoy
account with minimal funds, the same can also be made with
encryption volumes).

In general, assume a determined adversary could target you
specifically for your role: use confidential communication
to stay in touch with your team (so they know you're safe
daily), and establish a code word or protocol for
emergencies.

High-risk travelers might also arrange a ``check-in''
schedule with their security team or colleagues -- if you
don't check in by a certain time, they can take pre-agreed
actions (like disabling your accounts or alerting
authorities). This kind of planning is an extra safety net
when the stakes are especially high.

\section{Post-trip device rebuilding}

For highly targeted individuals, the safest course after
returning is to treat every device as compromised and
rebuild it, especially before reaching your safe zone
(home, work office). This involves wiping devices to
factory settings, flashing firmware if necessary, and
restoring data from known-good backups made before the
trip.

Consider using read-only operating systems or booting from
trusted live media during travel to reduce risk exposure.
Before your trip, you can create a cloned disk image of a
clean system state, so after traveling you can restore
your device to that exact low-level copy -- starting fresh
with a known secure baseline. This approach helps
eliminate stealthy malware or spyware that may have been
implanted.

Additionally, review system logs if available (security
apps or MDM solutions often report unusual access or
configuration changes during your absence).

As always, report any suspicious incidents promptly.
High-risk roles may require a debrief with your security
officer -- be transparent about any odd encounters or
possible security lapses to mitigate threats that could
have followed you home.
